\section{Measurement and Collapse}


This section organizes what measurement does to the conscious state, and in particular to qualia.

"Collapse" here does not directly mean the collapse of the wave function in quantum mechanics.

Collapse in this paper refers to the conversion of a live boundary state into a reportable, storable log.

That is, measurement is not an operation that finds consciousness, but an operation that transforms the consciousness opening state and sediments it into the slow layer.

\subsection{Measurement Requires Stabilization}


Measurement requires stabilization.

To measure something, it must be comparable, repeatable, and recordable.

Whether measuring neural activity, behavior, or subjective report, the object is stabilized in some form.

But the conscious state treated in this paper is not a completely stabilized object.

Consciousness is a temporal non-closure regime that opens when speed difference, coupling viscosity, internal delay, and meaning gradient enter a certain range.

This regime is not a static object but a dynamic relation.

Measurement is therefore not neutral observation.

Measurement is the operation of converting a moving relation into a recordable form.

\subsection{Measurement as Plasticity Closure}


The most important operation measurement performs is the closing of the plasticity gate.

Qualia appear at the boundary where fast-layer processing is sedimenting into the slow layer.

This boundary has not yet completely become a log.

But measurement converts that boundary into a reportable form.

Reporting. Choosing from options. Quantifying intensity. Recording as brain activity. Describing in language.

All of these are operations that sediment the live boundary state into the slow layer.

Measurement therefore does not extract qualia as they are.

Measurement logs qualia.

At that point, qualia disappear as qualia and a residue remains instead.

\subsection{From Qualia Live to Sedimented Residue}


This conversion can be expressed as:

\begin{equation}
\text{qualia live} \rightarrow \text{measurement} \rightarrow \text{sedimented residue}
\end{equation}

"Qualia live" is the live foreground appearing at the sedimentation frontier.

"Sedimented residue" is the log, report, label, number, or neural summary that remains after the boundary state has been sedimented by measurement into the slow layer.

The crucial point is that the residue is not meaningless.

The logs that remain after measurement can indeed be analyzed, compared, and subjected to statistical treatment.

But they are not qualia themselves.

They are the traces left after qualia have passed through.

What measurement can access is therefore not the substance of qualia but the passage traces of qualia.

\subsection{Why Precision Makes Qualia Disappear}


Increasing measurement precision can make qualia seem to disappear.

This is not because measurement is imperfect.

Rather, it is because measurement is succeeding.

Precise measurement stabilizes the object more strongly.

The stronger the stabilization, the more readily the live boundary sediments into the slow layer.

As a result, what measurement yields becomes a more definite log, a more definite label, a more definite number.

But that definiteness is not the definiteness of qualia themselves.

It is the definiteness of the residue after qualia have sedimented.

For this reason, measurement clarifies qualia while simultaneously erasing them as qualia.

Here lies the paradox of qualia measurement.

\subsection{Measurement and the Consciousness Window}


This measurement-induced transformation can also be understood from the perspective of the consciousness window.

A stable conscious state exists in the following intermediate region:

\begin{equation}
0 < |\Delta v| < \Delta v_{\max}, \quad \mu_{\min} < \mu < \mu_{\max}, \quad \tau_{\min} < \tau < \tau_{\max}, \quad \mathcal{T}_{\min} < \mathcal{T}_{\Phi} < \mathcal{T}_{\max}
\end{equation}

Measurement moves the live state within this window toward a slower, more stable log.

That is, measurement changes the relation among $\Delta v$, $\tau$, $\mu$, and $\mathcal{T}_{\Phi}$.

In particular, report and recording push the boundary state toward the slow layer.

As a result, the measurement target is converted from a live opening into a sedimented residue.

In this sense, measurement is not an operation that peers into consciousness from outside without disturbing it.

Measurement is an operation that changes the regime of the conscious state.

\subsection{Measurement Is Not Error}


One point of caution here.

The fact that measurement converts qualia into residue does not render measurement meaningless.

Measurement is not failure.

Measurement is the necessary transformation for handling live boundary states.

The problem is mistaking measurement results for qualia themselves.

Subjective reports, reaction times, neural activity, and linguistic descriptions obtained through measurement are all important.

But they are not boundary phenomena themselves; they are logs after boundary phenomena have sedimented.

The position of this paper is therefore not anti-measurement.

Rather, it is the position of precisely repositioning what measurement is measuring.

Measurement does not measure qualia.

Measurement measures the sedimented structure remaining after qualia have passed through.

\subsection{Introspection as Internal Measurement}


This structure applies not only to external measurement but also to introspection.

Introspection is the operation of observing, verbalizing, and grasping one's own conscious state.

But introspection is also measurement.

Directing attention to a sensation. Putting it into words. Evaluating its intensity. Considering its cause.

All of these operations sediment the live boundary state into the slow layer.

What introspection captures is therefore also not qualia themselves but the structure remaining after qualia have sedimented.

This does not mean introspection is unimportant.

Introspection is a powerful method of log updating in the slow layer.

But introspection does not store qualia; it logs them.

Confusing this point leads to the illusion that introspection directly sees the substance of consciousness.

\subsection{Experimental Implication}


In the framework of this paper, experimental design in consciousness research also takes on a slightly different perspective.

What should be measured is not qualia themselves.

What should be measured is under what conditions qualia appear, under what conditions they sediment, and what residue they leave.

The experimentally important quantities are therefore:

\begin{enumerate}
\item The delay until live experience is converted into a reportable log.
\item The transformation of experiential content by the strength of report demands.
\item Changes in residue due to differences in semantic load.
\item Changes in $\tau$ or $\mathcal{T}_{\Phi}$ due to attention manipulation.
\item Changes in neural activity patterns before and after report.
\item How experience is stabilized through verbalization.
\end{enumerate}

These are not attempts to capture qualia directly.

They are attempts to investigate the process by which qualia sediment.

\subsection{Relation to the Measurement Problem}


The collapse described in this section is not a direct interpretation of the quantum measurement problem.

But there is a structural similarity.

In both cases, measurement is not a neutral operation that preserves the object as it is.

Measurement is the operation of stabilizing a state and converting it into a recordable form.

In the case of qualia, this operation converts the live boundary state into sedimented residue.

The measurement problem in consciousness research is therefore not the problem of where qualia are hiding.

It is the problem of how boundary states are logged by measurement.

\subsection{Summary}


This section organized measurement not as external observation of the conscious state but as transformation of the opening state.

\begin{itemize}
\item Qualia are the live boundary state appearing at the sedimentation frontier.
\item Measurement sediments that boundary state into the slow layer.
\item What remains after measurement is sedimented residue, not qualia themselves.
\item As measurement precision increases, residue becomes more definite, but live qualia are closed.
\item Introspection is also internal measurement and logs qualia.
\item Measurement is not error but transformation for handling boundary states.
\item What should be experimentally treated is not qualia themselves but the conditions and residue of qualia sedimentation.
\end{itemize}

Measurement is therefore not the capture of consciousness as it is.

Measurement is the operation that closes the temporal non-closure opening and converts the live boundary into slow-layer logs.
